% Modello "Relazione" — LaTeX Editor Offline
% Contenuto segnaposto originale: libero da usare e modificare.
\documentclass[12pt,a4paper]{report}

\usepackage[T1]{fontenc}
\usepackage[utf8]{inputenc}
\usepackage[italian]{babel}
\usepackage{amsmath,amssymb}
\usepackage{graphicx}
\usepackage{booktabs}
\usepackage[hidelinks]{hyperref}

\title{Relazione di esempio}
\author{Nome Autore}
\date{\today}

\begin{document}

\maketitle

\begin{abstract}
  Sommario della relazione: obiettivo, contesto, attività svolta e
  risultati principali, in mezza pagina al massimo.
\end{abstract}

\tableofcontents

\chapter{Introduzione}
Descrivi lo scopo della relazione e il contesto in cui nasce il lavoro.
Una citazione di esempio: \cite{esempio2024}.

\chapter{Analisi}
\section{Requisiti}
Elenca i requisiti osservati.

\section{Vincoli}
Descrivi i vincoli tecnici o organizzativi.

\chapter{Attività svolta}
Descrivi le fasi del lavoro. Una figura si inserisce così:

\begin{figure}[htbp]
  \centering
  % \includegraphics[width=0.7\linewidth]{esempio.png}
  \fbox{\parbox[c][4cm][c]{0.7\linewidth}{\centering spazio per una figura}}
  \caption{Figura segnaposto.}
  \label{fig:segnaposto}
\end{figure}

Un riferimento alla figura: figura~\ref{fig:segnaposto}.

\chapter{Risultati}
Riporta i risultati, eventualmente con una tabella:

\begin{table}[htbp]
  \centering
  \caption{Risultati di esempio.}
  \begin{tabular}{lrr}
    \toprule
    Prova & Prima & Dopo \\
    \midrule
    A     & 10    & 12   \\
    B     & 20    & 25   \\
    \bottomrule
  \end{tabular}
\end{table}

\chapter{Conclusioni}
Chiudi la relazione con le osservazioni finali e i possibili sviluppi.

\bibliographystyle{plain}
\bibliography{bibliography}

\end{document}
